% Diagrama TikZ: Alinhamento OE → Q → H (6 colunas)
% Requer \usepackage{tikz} e \usetikzlibrary{...} no preâmbulo

\begin{tikzpicture}[
  % ── Estilos dos nós ──────────────────────────────────────────
  OGstyle/.style={
    draw=black!80, fill=black!8, rounded corners=3pt,
    font=\footnotesize\bfseries, text width=16.4cm, align=center,
    minimum height=1.1cm, inner sep=4pt},
  OEstyle/.style={
    draw=blue!70!black, fill=blue!8, rounded corners=2pt,
    font=\scriptsize, text width=2.25cm, align=center,
    minimum height=2.6cm, inner sep=2pt},
  Qstyle/.style={
    draw=orange!70!black, fill=orange!8, rounded corners=2pt,
    font=\scriptsize, text width=2.25cm, align=center,
    minimum height=2.6cm, inner sep=2pt},
  Hstyle/.style={
    draw=green!50!black, fill=green!8, rounded corners=2pt,
    font=\scriptsize, text width=2.25cm, align=center,
    minimum height=2.6cm, inner sep=2pt},
  CAPstyle/.style={
    draw=gray!60, fill=gray!10, rounded corners=2pt,
    font=\scriptsize\itshape, text width=2.25cm, align=center,
    minimum height=0.7cm, inner sep=2pt},
  CVstyle/.style={
    draw=black!80, fill=black!8, rounded corners=3pt,
    font=\footnotesize\bfseries, text width=16.4cm, align=center,
    minimum height=1.0cm, inner sep=4pt},
  % ── Setas ────────────────────────────────────────────────────
  arr/.style={-{Stealth[length=4pt]}, thick, draw=gray!70},
  arrconv/.style={-{Stealth[length=4pt]}, thick, draw=black!50, dashed},
  % ── Rótulos de coluna ───────────────────────────────────────
  coltitle/.style={font=\scriptsize\bfseries\scshape, text=gray!60!black},
]

% ══════════════════════════════════════════════════════════════
% OBJETIVO GERAL
% ══════════════════════════════════════════════════════════════
\node[OGstyle] (OG) {%
  \textsc{Objetivo Geral}\\[2pt]
  Desenvolver e validar um modelo integrado de decodifica\c{c}\~ao e salvaguarda, articulando\\
  aprendizado de m\'aquina, l\'ogica fuzzy e modelagem psicom\'etrica,\\
  para documentar e proteger ativos intang\'\i veis em SSAT quilombolas\\
  do Semi\'arido Nordeste\,II (Bahia)};

% ══════════════════════════════════════════════════════════════
% RÓTULOS DE COLUNA
% ══════════════════════════════════════════════════════════════
\node[coltitle, below=0.55cm of OG.south west, anchor=west, xshift=0.1cm]
  (labOE) {Objetivos Espec\'\i ficos};

% ══════════════════════════════════════════════════════════════
% LINHA 1 — OBJETIVOS ESPECÍFICOS
% ══════════════════════════════════════════════════════════════
\node[OEstyle, below=1.0cm of OG.south west, anchor=north west, xshift=0.0cm]
  (OE1) {\textbf{OE\,1}\\[1pt]
  Mapear estado da arte de ML em SAT (PRISMA-ScR), lacunas em FAIR, XAI e valida\c{c}\~ao};

\node[OEstyle, right=0.20cm of OE1] (OE2)
  {\textbf{OE\,2}\\[1pt]
  Investigar como dimens\~oes de vulnerabilidade se associam \`a degrada\c{c}\~ao de SSAT (PRISMA\,2020)};

\node[OEstyle, right=0.20cm of OE2] (OE3)
  {\textbf{OE\,3}\\[1pt]
  Adaptar WOCAT-SLM e elicitar vari\'aveis via Delphi para decodifica\c{c}\~ao computacional};

\node[OEstyle, right=0.20cm of OE3] (OE4)
  {\textbf{OE\,4}\\[1pt]
  Desenvolver motor fuzzy + TRI gerando o \'Indice de Salvaguarda Biocultural (ISB)};

\node[OEstyle, right=0.20cm of OE4] (OE5)
  {\textbf{OE\,5}\\[1pt]
  Avaliar efic\'acia probat\'oria e aceita\c{c}\~ao institucional dos dossi\^es regulat\'orios};

\node[OEstyle, right=0.20cm of OE5] (OE6)
  {\textbf{OE\,6}\\[1pt]
  Validar ado\c{c}\~ao territorial via design participativo com comunidades quilombolas};

% ══════════════════════════════════════════════════════════════
% LINHA 2 — QUESTÕES DE PESQUISA
% ══════════════════════════════════════════════════════════════
\node[Qstyle, below=0.55cm of OE1] (Q1)
  {\textbf{Q\,1}\\[1pt]
  Em que medida ML apresenta prontid\~ao operacional para SAT?};

\node[Qstyle, below=0.55cm of OE2] (Q2)
  {\textbf{Q\,2}\\[1pt]
  Como as dimens\~oes $V_1$--$V_6$ se associam \`a degrada\c{c}\~ao e quais fatores as modulam?};

\node[Qstyle, below=0.55cm of OE3] (Q3)
  {\textbf{Q\,3}\\[1pt]
  Como WOCAT-SLM + Delphi convertem saberes t\'acitos em vari\'aveis computacionais?};

\node[Qstyle, below=0.55cm of OE4] (Q4)
  {\textbf{Q\,4}\\[1pt]
  O ISB fuzzy produz maior ader\^encia comunit\'aria que classifica\c{c}\~oes lineares?};

\node[Qstyle, below=0.55cm of OE5] (Q5)
  {\textbf{Q\,5}\\[1pt]
  Os dossi\^es atendem requisitos probat\'orios de inst\^ancias regulat\'orias?};

\node[Qstyle, below=0.55cm of OE6] (Q6)
  {\textbf{Q\,6}\\[1pt]
  Como \textit{participatory design} afeta apropria\c{c}\~ao e ado\c{c}\~ao continuada?};

% ══════════════════════════════════════════════════════════════
% LINHA 3 — HIPÓTESES
% ══════════════════════════════════════════════════════════════
\node[Hstyle, below=0.55cm of Q1] (H1)
  {\textbf{H\,1}\\[1pt]
  Prontid\~ao ML insuficiente por lacunas em governan\c{c}a de dados dos estudos prim\'arios};

\node[Hstyle, below=0.55cm of Q2] (H2)
  {\textbf{H\,2}\\[1pt]
  Ao menos uma dimens\~ao $V_1$--$V_6$ difere significativamente das demais em magnitude de efeito};

\node[Hstyle, below=0.55cm of Q3] (H3)
  {\textbf{H\,3}\\[1pt]
  WOCAT-SLM + Delphi produz consenso superior a m\'etodos n\~ao estruturados};

\node[Hstyle, below=0.55cm of Q4] (H4)
  {\textbf{H\,4}\\[1pt]
  ISB fuzzy apresenta maior concord\^ancia comunit\'aria que classifica\c{c}\~oes lineares};

\node[Hstyle, below=0.55cm of Q5] (H5)
  {\textbf{H\,5}\\[1pt]
  Modelo eleva qualidade t\'ecnica (padroniza\c{c}\~ao, completude, rastreabilidade) vs.\ convencional};

\node[Hstyle, below=0.55cm of Q6] (H6)
  {\textbf{H\,6}\\[1pt]
  Design participativo eleva ado\c{c}\~ao continuada vs.\ ferramentas ex\'ogenas};

% ══════════════════════════════════════════════════════════════
% LINHA 4 — CAPÍTULOS (ARTIGOS)
% ══════════════════════════════════════════════════════════════
\node[CAPstyle, below=0.45cm of H1] (C1) {Cap\'\i tulo 1};
\node[CAPstyle, below=0.45cm of H2] (C2) {Cap\'\i tulo 2};
\node[CAPstyle, below=0.45cm of H3] (C3) {Cap\'\i tulo 3};
\node[CAPstyle, below=0.45cm of H4] (C4) {Cap\'\i tulo 4};
\node[CAPstyle, below=0.45cm of H5] (C5) {Cap\'\i tulo 5};
\node[CAPstyle, below=0.45cm of H6] (C6) {Cap\'\i tulo 6};

% ══════════════════════════════════════════════════════════════
% CONVERGÊNCIA
% ══════════════════════════════════════════════════════════════
\node[CVstyle, below=0.9cm of C3.south east] (CV) {%
  \textsc{Valida\c{c}\~ao Convergente}\\[1pt]
  Modelo integrado: prontid\~ao ML\,(H1) $\rightarrow$
  dimens\~oes $V_1$--$V_6$\,(H2) $\rightarrow$
  vari\'aveis Delphi\,(H3) $\rightarrow$ ISB fuzzy\,(H4) $\rightarrow$
  qualidade t\'ecnica documental\,(H5) $\rightarrow$ ado\c{c}\~ao territorial\,(H6)};

% ══════════════════════════════════════════════════════════════
% SETAS VERTICAIS
% ══════════════════════════════════════════════════════════════
\foreach \i in {1,...,6}{
  \draw[arr] (OG.south) -- ++(0,-0.35) -| (OE\i.north);
  \draw[arr] (OE\i.south) -- (Q\i.north);
  \draw[arr] (Q\i.south) -- (H\i.north);
  \draw[arr] (H\i.south) -- (C\i.north);
  \draw[arrconv] (C\i.south) -- ++(0,-0.25) -| (CV.north);
}

% ══════════════════════════════════════════════════════════════
% RÓTULOS LATERAIS
% ══════════════════════════════════════════════════════════════
\begin{scope}[coltitle]
  \node[rotate=90, anchor=south] at ($(OE1.west)+(-0.45,0)$) {Objetivos};
  \node[rotate=90, anchor=south] at ($(Q1.west)+(-0.45,0)$)  {Quest\~oes};
  \node[rotate=90, anchor=south] at ($(H1.west)+(-0.45,0)$)  {Hip\'oteses};
  \node[rotate=90, anchor=south] at ($(C1.west)+(-0.45,0)$)  {Artigos};
\end{scope}

\end{tikzpicture}
